\section{Formulación del problema}

\subsection{Problema general}
¿Qué configuración de una red neuronal convolucional ligera ---combinación de
arquitectura (YOLOv8n o YOLOv11) y esquema de cuantización (FP32 o INT8)--- ofrece
el mejor compromiso entre precisión (incluida la robustez), eficiencia
computacional y oportunidad de la retroalimentación docente en dispositivos
móviles?

\subsection{Problemas específicos}
\begin{enumerate}
  \item ¿Qué configuración ofrece mayor precisión y robustez, medida con
  mAP@0.5, F1 macro, Kappa de Cohen y AUC-ROC macro, en condiciones normales y
  adversas?
  \item ¿En qué medida la cuantización INT8 reduce la latencia y el consumo de
  memoria sin degradar la precisión más allá de un umbral aceptable?
  \item ¿Qué efecto tiene el sistema en la oportunidad de la retroalimentación
  docente (tiempo de revisión por página) en aulas inclusivas de Moquegua e Ilo?
\end{enumerate}
